\documentclass{exam}

\usepackage{Haust2016verkefnablöð}

\title{Stærðfræðimynstur í tölvunarfræði \\ Skilaverkefni 7}
\author{}

\printanswers

\begin{document}
\maketitle
\thispagestyle{empty} 

Skila skal þessu verkefni á vefnum \href{https://gradescope.com/}{Gradescope}. Aðgangskóði fyrir námskeiðið er \textbf{926WD9}.


\section{Spurningar}

\begin{questions}

\subsection{Kafli 6.1}

\question 
\begin{enumerate}[a)]
 \item Hversu mörg föll eru til frá mengi með 10 stökum til mengis með 4 stökum?
 \item Gefið er mengi með 10 stökum. Hversu mörg af hlutmengjum þess hafa meira en eitt stak?
\end{enumerate}

\paragraph{Í bók:} Dæmi 6.1.34c, 6.1.40

\subsection{Kafli 6.2}

\question Gefnar eru 51 heiltala, hver þeirra er á lokaða bilinu frá 0 til 99. Sýnið að a.m.k. tvær talnanna séu aðliggjandi, þ.e.a.s. að til sé $n$ þannig að bæði $n$ og $n+1$ séu meðal gefnu talnanna.

\paragraph{Í bók:} Byggt á dæmi 6.2.44

\subsection{Kafli 6.3}

\question Hversu margir bitastrengir af lengd 19 innihalda nákvæmlega fimm 0 og fjórtán 1 ef á eftir hverju 0 þurfa að koma tvö 1?

\paragraph{Í bók:} Dæmi 6.3.36

\newpage
\subsection{Kafli 6.4}

\question Sýnið að

\[
 \binom{2n}{2} = 2\binom{n}{2} + n^2
\]

\begin{enumerate}[a)]
 \item með bókstafareikningi
 \item með tvöfaldri talningu
\end{enumerate}

\paragraph{Í bók:} Dæmi 6.4.28

\subsection{Kafli 8.2}

Finnið lokaða formúlu fyrir $a_n$ í rakningarvenslunum
\[
 a_n = a_{n-1} + 6a_{n-2}
\]
þegar $n \geq 2$. Upphafsskilyrði eru $a_0 = 3$ og $a_1 = 6$.

\paragraph{Í bók:} Dæmi 8.2.4a

\subsection{Kafli 8.3}

Notið ``master theorem'' til að finna stóra-O mat á $f(n) = 9f\left(\frac{n}{3}\right) + n$.

\paragraph{Í bók:} Þetta dæmi er ekki í bókinni.
\end{questions}

\vfill
\includegraphics[width=0.5\linewidth]{hi-von-logo}

\end{document}